% cat-location-lattice.tex — the funnel from "where code starts" to "where it ships".
% Rendered with the tikz MCP server (compile_tikz); the PNG drops in assets/ as
% cat-location-lattice.png.
%
% INK: white on a transparent background — the site's theme is dark (--maincolor
% #000000, panels #191C24). Edge labels are UNFILLED and offset clear of the strokes.
%
% THE FIGURE, IN ONE PICTURE. The location side of a distribution rewrite, read the
% way the type-part/value-residue split is read on the code side: a DECLARED test
% (admissible — some authority's scheme has a dialer) and a REALIZED test (serves —
% the dialed peer's client overrides every instruction the placement requires, and is
% online). origin is where the code starts (this machine's listeners); reachable is
% where it can go (roster ∪ dialed peers); shardTargets is the subset a plan actually
% ships to, addressed by WIRE authority (never the local map key).
%
%       origin              where the code STARTS      Places.self()
%         │                    (this machine's listeners)
%         ▼
%       reachable            where it CAN GO           Places.reachable()
%         │                    roster ∪ dialed peers
%         │
%         ├─ admissible       DECLARED: some authority's scheme resolves a dialer
%         ▼
%       serves               REALIZED: the client overrides every placement.inst,
%         │                     and the link is online          Places.serves()
%         ▼
%       shardTargets         dialed (has host) + serves, shipped to the WIRE authority

\begin{tikzpicture}[
  c/.style={draw, rounded corners=2pt, thick, minimum width=26mm, minimum height=6.5mm,
            font=\small, inner sep=1.2mm},
  e/.style={-{Stealth[length=2.2mm]}, thick},
  el/.style={font=\scriptsize, inner sep=1pt},
  sub/.style={font=\tiny, text=white!60!black, align=left, inner sep=1pt},
  note/.style={font=\scriptsize, inner sep=1pt},
  color=white
]

% ---- the funnel ----
\node[c] (o)  at (0,0)  {\texttt{origin}};
\node[c] (r)  at (0,-1.6) {\texttt{reachable}};
\node[c] (s)  at (0,-3.2) {\texttt{serves}};
\node[c] (t)  at (0,-4.8) {\texttt{shardTargets}};

\draw[e] (o) -- node[el,right=2pt]{$\mathrm{self}()$} (r);
\draw[e] (r) -- node[el,right=2pt]{$\mathrm{admissible}$} (s);
\draw[e] (s) -- node[el,right=2pt]{} (t);

% ---- right column: what each is ----
\node[sub, anchor=west, align=left] at (1.9,0)  {where the code STARTS\\ (this machine's listeners)};
\node[sub, anchor=west, align=left] at (1.9,-1.6) {where it CAN GO\\ (roster $\cup$ dialed peers)};
\node[sub, anchor=west, align=left] at (1.9,-3.2) {client overrides every\\ \texttt{placement.inst}, and is online};
\node[sub, anchor=west, align=left] at (1.9,-4.8) {dialed (has \texttt{host}) + serves,\\ shipped to the WIRE authority};

% ---- the declared / realized split ----
\node[sub, anchor=east, align=right] at (-1.9,-2.4) {DECLARED:\\ some scheme\\ resolves a dialer};
\node[sub, anchor=east, align=right] at (-1.9,-4.0) {REALIZED:\\ overrides + online};
\end{tikzpicture}
